%======================================================================
\chapter[Integration by parts: the Gauss--Ostrogradsky family]{Integration by parts:\\ the Gauss--Ostrogradsky family}\label{app:ibp}
%======================================================================
% Sources: "Why variational formulations matter", Section 11; notes C1.

Every weak formulation and every adjoint computation in this book rests on a
small family of classical identities. Throughout, $\Omega\subset\R^3$ is a
bounded domain with Lipschitz boundary $\partial\Omega$, outward unit normal
$\vn$ and surface measure $\dd s$~\cite{MarsdenTromba2011,Evans2010B}. Indices
refer to Cartesian components, with summation over repeated indices and
$(\cdot)_{,j}=\partial(\cdot)/\partial x_j$.

%======================================================================
\section{The divergence theorem}
\begin{theorem}[Gauss--Ostrogradsky]\label{thm:div}
For a vector field $\vF\in C^1(\overline\Omega;\R^3)$,
\[
  \int_\Omega\dive\vF\dd x=\int_{\partial\Omega}\vF\cdot\vn\dd s,
  \qquad\text{i.e.}\qquad
  \int_\Omega F_{i,i}\dd x=\int_{\partial\Omega}F_in_i\dd s .
\]
\end{theorem}
This is the master identity; every formula below follows from it and a
product rule. By density it extends to Sobolev fields, with boundary values
understood as traces.\footnote{On the names: the divergence theorem is due to
Gauss (1813, a special case) and Ostrogradsky (1826, the general proof),
independently of the curl theorem now called after Stokes. That theorem was
itself communicated by Kelvin to Stokes in a letter of 2 July 1850 and owes
its name to a question Stokes set in the 1854 Smith's Prize
examination~\cite{Katz1979}, hence the name Kelvin--Stokes. All of these are
special cases of the generalized Stokes theorem for differential forms,
$\int_Md\omega=\int_{\partial M}\omega$, which is why the divergence theorem
is sometimes called ``Stokes' theorem'' too. We stay with classical vector
calculus and keep the names separate, following~\cite{MarsdenTromba2011}.}

%======================================================================
\section{Scalar fields}
\begin{theorem}[Integration by parts, scalar times vector]\label{thm:ibp}
For $u\in C^1(\overline\Omega)$ and $\vF\in C^1(\overline\Omega;\R^3)$,
applying Theorem~\ref{thm:div} to $\dive(u\vF)=u\dive\vF+\nabla u\cdot\vF$,
\[
  \int_\Omega u\dive\vF\dd x=-\int_\Omega\nabla u\cdot\vF\dd x
  +\int_{\partial\Omega}u\,(\vF\cdot\vn)\dd s .
\]
\end{theorem}
With $\vF=\nabla u/\sqrt{1+|\nabla u|^2}$ it gives the minimal surface
equation (Section~\ref{sec:var-phys}); with $\vF=\vp$ the adjoint
$\Lambda^\ast\vp=-\dive\vp$ of Section~\ref{sec:var-dual}.

\begin{corollary}[Componentwise]\label{cor:comp}
For $i=1,2,3$ and $u,v\in C^1(\overline\Omega)$,
$\displaystyle\int_\Omega u\,v_{,i}\dd x=-\int_\Omega u_{,i}\,v\dd x+\int_{\partial\Omega}uv\,n_i\dd s$.
\end{corollary}

\begin{theorem}[Green's first identity]\label{thm:green1}
Taking $\vF=\nabla v$ in Theorem~\ref{thm:ibp},
\[
  \int_\Omega\nabla u\cdot\nabla v\dd x=-\int_\Omega u\,\Delta v\dd x
  +\int_{\partial\Omega}u\,\partial_nv\dd s .
\]
\end{theorem}
This is the identity that passes from the strong to the weak form of
$-\Delta u=f$ (Section~\ref{sec:var-weak}).

\begin{theorem}[Green's second identity]\label{thm:green2}
Subtracting Theorem~\ref{thm:green1} from itself with $u$ and $v$ exchanged,
\[
  \int_\Omega\bigl(u\,\Delta v-v\,\Delta u\bigr)\dd x
  =\int_{\partial\Omega}\bigl(u\,\partial_nv-v\,\partial_nu\bigr)\dd s .
\]
\end{theorem}
This symmetric form makes the Laplacian self-adjoint
(Exercise~\ref{exo:selfadj}). It is the continuum form of the reciprocity
theorem (Section~\ref{ssec:reciprocity}) and the starting point of the
reciprocity gap (Chapter~\ref{ch:rg}).

\begin{theorem}[Matrix coefficient field]\label{thm:ibp-matrix}
For a second-order tensor field $\tA\in C^1(\overline\Omega;\R^{3\times3})$ and
$u,v\in C^1(\overline\Omega)$, Theorem~\ref{thm:ibp} with $\vF=\tA\nabla u$ gives
\[
  \int_\Omega\tA\nabla u\cdot\nabla v\dd x
  =-\int_\Omega\dive(\tA\nabla u)\,v\dd x+\int_{\partial\Omega}v\,(\tA\nabla u)\cdot\vn\dd s .
\]
If $\tA$ is symmetric, exchanging $u$ and $v$ and subtracting gives the
second identity
$\int_\Omega\bigl(u\dive(\tA\nabla v)-v\dive(\tA\nabla u)\bigr)
=\int_{\partial\Omega}\bigl(u\,(\tA\nabla v)\cdot\vn-v\,(\tA\nabla u)\cdot\vn\bigr)$.
\end{theorem}
It is used for the general elliptic operator (Section~\ref{sec:var-elliptic})
and for the diffusion problem (Section~\ref{sec:diff}).

%======================================================================
\section{Tensor fields}
For a second-order tensor field $\sig$, the divergence is the vector
$(\dive\sig)_i=\sigma_{ij,j}$, and $\sig\vn$ has components $\sigma_{ij}n_j$.

\begin{theorem}[Tensor times vector]\label{thm:ibp-tensor}
For $\sig\in C^1(\overline\Omega;\R^{3\times3})$ and
$\vv\in C^1(\overline\Omega;\R^3)$,
\[
  \int_\Omega\sig:\nabla\vv\dd x=-\int_\Omega(\dive\sig)\cdot\vv\dd x
  +\int_{\partial\Omega}(\sig\vn)\cdot\vv\dd s .
\]
If $\sig$ is symmetric, $\sig:\nabla\vv=\sig:\eps[\vv]$ with
$\eps[\vv]=\tfrac12(\nabla\vv+\nabla\vv\T)$.
\end{theorem}
\begin{proof}
Apply Theorem~\ref{thm:div} to the vector $F_j=\sigma_{ij}v_i$:
$F_{j,j}=\sigma_{ij,j}v_i+\sigma_{ij}v_{i,j}$ and $F_jn_j=(\sigma_{ij}n_j)v_i$.
For symmetric $\sig$, the antisymmetric part of $\nabla\vv$ does not
contribute to $\sig:\nabla\vv$.
\end{proof}
This identity gives the principle of virtual work (Section~\ref{sec:el-vw}):
the divergence on symmetric tensors is minus the adjoint of the symmetric
gradient. With $\sig=\bbC:\eps[\vu]$ and $\bbC$ having the major symmetry, it
gives the Betti reciprocity theorem of elasticity,
\[
  \int_\Omega\bigl(\vu\cdot\dive\sig[\vv]-\vv\cdot\dive\sig[\vu]\bigr)\dd x
  =\int_{\partial\Omega}\bigl(\vu\cdot\sig[\vv]\vn-\vv\cdot\sig[\vu]\vn\bigr)\dd s .
\]

%======================================================================
\section{Domains with a cut}
When $\Omega$ contains a crack $\Gamma$ with lips $\Gamma^\pm$, the identities
hold in $\OG=\Omega\setminus\Gamma$, whose boundary consists of $\partial\Omega$
and of both lips. If $\vn$ is the normal of $\Gamma$ pointing towards
$\Gamma^+$, the outward normal of $\OG$ is $-\vn$ on $\Gamma^+$ and $+\vn$ on
$\Gamma^-$. For a field $\vF$ smooth on each side,
\[
  \int_{\OG}\dive\vF\dd x=\int_{\partial\Omega}\vF\cdot\vn\dd s
  -\int_\Gamma\jump{\vF}\cdot\vn\dd s,\qquad\jump{\vF}=\vF^+-\vF^- .
\]
This is how the jump of the potential across the crack appears in the
reciprocity gap identity~\eqref{eq:rg-rg}.

%======================================================================
\section{Where each identity is used}
\begin{center}
\renewcommand{\arraystretch}{1.15}
\begin{tabular}{ll}
\toprule
\sffamily where used & \sffamily identity\\
\midrule
weak form of $-\Delta u=f$ (Section~\ref{sec:var-weak}) & Green's first identity, Theorem~\ref{thm:green1}\\
symmetry of $a(u,v)$, reciprocity (Sections~\ref{sec:var-lm}, \ref{ssec:reciprocity}) & Green's second identity, Theorem~\ref{thm:green2}\\
minimal surface equation (Section~\ref{sec:var-phys}) & Theorem~\ref{thm:ibp}\\
general elliptic operator, diffusion (Sections~\ref{sec:var-elliptic}, \ref{sec:diff}) & Theorem~\ref{thm:ibp-matrix}\\
dual problem, $\Lambda^\ast=-\dive$ (Section~\ref{sec:var-dual}) & Theorem~\ref{thm:ibp}\\
virtual work, Navier equations (Sections~\ref{sec:el-vw}, \ref{sec:el-energy}) & Theorem~\ref{thm:ibp-tensor}\\
reciprocity gap (Section~\ref{sec:rg-gap}) & Green's second identity in $\OG$\\
\bottomrule
\end{tabular}
\end{center}

\paragraph{The classical trio.} The divergence theorem is the
three-dimensional member of a family that also contains the \emph{gradient
theorem} (fundamental theorem of line integrals) and \emph{Stokes' theorem}
(the flux of $\operatorname{curl}\vF$ through a surface equals the
circulation of $\vF$ along its boundary curve). Together they generalize the
fundamental theorem of calculus to higher dimensions~\cite{MarsdenTromba2011}.
Only the divergence theorem is needed here, since every variational problem
of the book is built from $\nabla$ and its adjoint $-\dive$.

%======================================================================
\begin{chapterbib}{9}
\bibitem{Evans2010B} L.~C.~Evans, \emph{Partial Differential Equations}, 2nd ed., AMS, 2010.
\bibitem{Katz1979} V.~J.~Katz, The history of Stokes' theorem, \emph{Math. Mag.} 52 (1979) 146--156.
\bibitem{MarsdenTromba2011} J.~E.~Marsden, A.~J.~Tromba, \emph{Vector Calculus}, 6th ed., W.~H.~Freeman, 2011.
\end{chapterbib}
